\chapter{研究设计、估计方法与不确定性}
\label{ch:design-estimation}

\begin{learninggoals}
读完本章，你应当能够：从已经定义并识别的目标量选择研究设计和估计器；比较回归、匹配、加权、双重稳健与机器学习辅助估计；报告抽样、模型和结构不确定性。
\end{learninggoals}

\section{从识别到估计：设计、模型与不确定性}

\subsection{结果建模、匹配与倾向得分}

当标准化公式
\[
E\{Y(a)\}=\int E(Y\mid A=a,L=l)\,dP(l)
\]
已经由一致性、可交换性和正值性识别后，我们仍有多种估计方式。它们使用同一识别逻辑，却对结果模型、处理模型、平滑性和有限样本表现作不同选择。比较方法时首先应确认它们估计同一目标人群、同一效应尺度和同一处理策略。

“机器学习方法”“因果森林”或“倾向得分匹配”不是目标量。一个复杂估计器可以非常准确地估计错误或含混的目标。反之，设计清楚的随机试验用简单均值差就可能给出高质量答案。

\paragraph{结果回归与g计算}

结果回归估计条件均值
\[
m_a(l)=E(Y\mid A=a,L=l),
\]
然后对目标人群的协变量分布标准化：
\[
\widehat\mu_a=\frac{1}{n}\sum_{i=1}^n \widehat m_a(L_i),
\qquad
\widehat{ATE}=\widehat\mu_1-\widehat\mu_0.
\]
这称为g计算或回归标准化。它直接模拟每个样本单位在两个处理状态下的预测结果，再求人群平均。

若 $m_a(l)$ 模型正确，估计可以一致。传统线性回归容易把模型系数误作效应：存在非线性、交互或非可合并链接时，处理系数未必等于边缘ATE。g计算通过显式标准化使目标更清楚，也允许使用灵活学习器。
模型诊断不能只看整体预测误差。因果估计需要在每个处理状态下预测目标人群潜在结果；缺少重叠的区域正是反事实外推最强之处，即使总体预测良好也可能偏差很大。

\paragraph{匹配与设计阶段的可比性}

匹配为每个处理单位寻找协变量相似的对照，试图重建类似随机试验的比较。可按精确变量、Mahalanobis距离或倾向得分匹配。设计阶段通常不使用结果数据来反复挑选匹配规则，减少结果导向的分析选择。
匹配后必须检查协变量平衡、共同支持、丢弃单位和目标人群变化。标准化均值差比逐变量显著性检验更适合平衡诊断，因为显著性受样本量影响。若大量处理单位找不到可比对照，得到的效应可能只针对重叠子人群，而非原定ATT或ATE。
最近邻匹配并不自动无偏。连续协变量的近似匹配存在残差差异，匹配后分析还需考虑成对或聚类结构，带放回匹配和一对多匹配改变权重。把匹配索引后简单做普通独立样本检验，可能给出错误标准误。

\paragraph{倾向得分与加权}

\term{倾向得分}是
\[
e(L)=P(A=1\mid L).
\]
在条件可交换性成立时，给定真实倾向得分可使已测协变量在处理组间平衡\parencite{rosenbaumrubin1983}。它压缩的是处理赋值相关信息，不是结果风险；倾向得分相近不保证未测混杂平衡。

ATE的逆概率处理权重可写为
\[
w_i=\frac{A_i}{e(L_i)}+\frac{1-A_i}{1-e(L_i)}.
\]
加权后构造一个处理与已测 $L$ 独立的伪人群，再比较加权结果均值。ATT、ATC和重叠人群使用不同权重；权重公式必须与目标一致。

倾向模型的目标是平衡，而不是最大化处理预测准确率。一个几乎完美预测处理的模型可能产生极端权重，暴露正值性问题。应报告倾向分布、权重分位数、有效样本量和加权后平衡。稳定化和截断可改善方差，但必须说明它们对目标与偏差的影响。匹配估计的渐近性质与倾向方法的应用诊断分别可见\textcite{abadieimbens2006}和\textcite{austin2011}。

\subsection{双重稳健、正交化与试验估计}

增广逆概率加权（AIPW）结合结果模型 $m_a(L)$ 和倾向模型 $e(L)$。处理水平 $a$ 的均值可估为
\[
\widehat\mu_a=\frac1n\sum_{i=1}^n
\left[
\widehat m_a(L_i)+
\frac{\mathbb I(A_i=a)}{\widehat P(A_i=a\mid L_i)}
\{Y_i-\widehat m_a(L_i)\}
\right].
\]

若结果模型或处理模型至少一个正确，在常规条件下估计仍可一致，因此称为\term{双重稳健}。这不是“两个错误模型相互抵消”的保证，也不抵抗未测混杂、处理定义错误或同时严重错设。有限样本中极端权重仍可能使估计不稳定\parencite{bangrobins2005}。
目标最大似然估计（TMLE）以初始结果模型为起点，用倾向信息沿目标量相关方向更新，使最终估计满足有效影响函数方程\parencite{vanderlaanrose2011}。它同样依赖识别假设，但在使用灵活学习器时提供系统的目标化估计与推断框架。

\paragraph{Double Machine Learning}

高维协变量下，直接把机器学习预测带入因果参数公式可能产生正则化偏差，并因同一数据既训练又评估而过拟合。\term{Double/Debiased Machine Learning}使用\term{Neyman正交}得分，使目标估计对干扰函数的小误差一阶不敏感，并结合\term{交叉拟合}：在一部分样本训练处理和结果模型，在未参与训练的部分计算得分，再轮换合并\parencite{chernozhukov2018}。

以部分线性模型
\[
Y=\theta A+g(L)+\varepsilon,\qquad E(\varepsilon\mid A,L)=0,
\]
\[
A=m(L)+v,\qquad E(v\mid L)=0
\]
为例，Robinson残差化并不是先直接估计结构函数 $g(L)$ 再从 $Y$ 中减去它。应定义两个可由观测分布估计的干扰函数
\[
\ell(L)=E(Y\mid L),\qquad m(L)=E(A\mid L).
\]
由于 $\ell(L)=\theta m(L)+g(L)$，有
\[
Y-\ell(L)=\theta\{A-m(L)\}+\varepsilon.
\]
DML在训练折用机器学习估计 $\ell$ 与 $m$，再在留出折用残差化后的 $Y-\widehat\ell(L)$ 与 $A-\widehat m(L)$ 构造正交得分并估计 $\theta$。正交性降低干扰函数误差传到 $\theta$ 的一阶影响；交叉拟合放宽对学习器复杂度的经验过程限制。

DML不是“先用ML去混杂就得到因果”。它仍需要可交换性、重叠和正确得分结构；不同DML模型对应不同参数。若结构不是部分线性，$\theta$ 的解释需要另行说明。

\paragraph{随机试验中的估计细节}

随机化识别简单均值差，但协变量调整可提高精度。应预先指定调整变量和模型，使用对随机化设计有效的标准误。分层、整群或不等概率随机化需要在分析中反映设计；回归调整的稳健性仍需依照随机化结构解释\parencite{lin2013}。
依从性不足时，\term{意向治疗效应}比较分配策略，保留随机化优势。按实际治疗比较会重新引入选择。若随机分配只通过实际治疗影响结果，且满足单调性等条件，可把分配作为工具变量估计服从者平均因果效应。这个目标与意向治疗效应不同，不应因后者“被稀释”就自动替换。
多重结局、亚组和中期分析增加选择性报告风险。预注册、分析计划和对探索性结果的清楚标记，是估计可信度的一部分。

\subsection{方法选择、数值诊断与分层不确定性}

置信区间通常反映在模型和识别假设成立时的抽样变异。它不包括未测混杂、测量误差、处理版本不明确、图结构错误和模型选择空间。一个极窄区间可能精确地围绕有偏目标。

完整的不确定性报告至少包括：

\begin{itemize}
  \item 点估计与决策相关的绝对、相对尺度；
  \item 抽样置信区间或兼容区间；
  \item 重叠、权重和协变量平衡诊断；
  \item 替代合理模型或调整集结果；
  \item 未测混杂、失访与测量误差敏感性分析；
  \item 从样本到目标人群的适用限制。
\end{itemize}

$p$值只是在特定模型下衡量数据与零假设的不相容程度，不说明效应大小、因果假设可信度或政策价值。统计不显著不等于无效，显著也不等于重要。

\paragraph{一张方法选择表}

\begin{center}
\small
\begin{tabularx}{\textwidth}{>{\bfseries}p{2.2cm}X X}
\toprule
方法 & 主要优势 & 主要风险\\
\midrule
g计算 & 直接对应标准化目标，易表达非线性与交互 & 结果模型外推和错设\\
匹配 & 强调设计与共同支持，诊断直观 & 丢弃样本、残差不平衡、推断处理复杂\\
IPW & 构造目标伪人群，适合纵向扩展 & 极端权重和倾向模型错设\\
AIPW/TMLE & 双重稳健，可结合灵活学习 & 两模型同时失败、有限样本不稳定\\
DML & 高维干扰函数、正交得分和交叉拟合 & 目标模型仍需正确，不能修复未测混杂\\
\bottomrule
\end{tabularx}
\end{center}

\paragraph{数值例子：权重为何会失控}

假设一个样本有1000人，其中某高风险层只有20人，19人接受治疗、1人未治疗。若估计ATE，未治疗者权重约为 $1/(1-0.95)=20$，一个人的结局会代表该层大量反事实信息。若其结局恰为异常值，估计剧烈变化。
这不是简单“模型方差太大”，而是数据几乎没有回答“高风险者不治疗会怎样”。截断权重到10可稳定计算，却相当于减少该人的代表性，并引入偏差。把目标改为重叠人群效应可能更诚实：它关注各处理都有真实选择机会的人，而不外推到几乎确定治疗者。

权重诊断应在看结果前完成。报告最大值、中位数、尾部分位数和有效样本量
\[
n_{\mathrm{eff}}=\frac{(\sum_iw_i)^2}{\sum_iw_i^2}.
\]
原始样本1000人若有效样本量只有120，精度不应按1000人的直觉解释。

\subsection{稳健性、标准误与决策解释}

AIPW中的结果预测为每个人填充两个处理下的条件均值，权重残差项校正结果模型在实际处理组中的系统误差。若倾向模型正确，残差加权可修正错误结果模型；若结果模型正确，残差平均为零，即使倾向模型错也可保持一致。
但“至少一个正确”是在总体与规则条件下的渐近命题。有限样本里，一个严重错误结果模型与极端错误倾向模型可以共同产生灾难结果；灵活学习器若不交叉拟合，会因过拟合破坏标准推断。双重稳健也完全不涉及未测混杂和一致性。
实践中应同时诊断两模型：倾向模型看平衡与重叠，结果模型看分层校准与外推；再比较g计算、IPW和AIPW。三者分歧大不是挑选最喜欢结果的理由，而是模型依赖警报。

\paragraph{标准误应尊重数据生成层次}

班级内学生共享教师、同伴和处理实施，观测不独立。即使处理在学生层分配，结果误差也可能聚类；若处理在学校层随机，独立信息单位更接近学校数量。把几千名学生当独立可大幅低估标准误。
重复测量、匹配组、网络和多阶段抽样都有类似问题。标准误、bootstrap和交叉验证划分必须保持依赖单元。例如同一患者的多次记录不能随机分到训练与验证折，否则泄漏个体背景并夸大泛化。
小簇数量时，普通簇稳健标准误可能仍偏小，需要小样本修正、随机化推断或适合层级的模型。异方差稳健方差的小样本修正与聚类稳健推断的适用条件分别见\textcite{mackinnonwhite1985}和\textcite{cameronmiller2015}。设计信息不是采集完后可丢弃的元数据，而是推断分布的一部分。

\paragraph{从估计结果到决策}

效应估计需与基线风险、成本和损害结合。风险比0.8在基线风险50\%时对应10个百分点绝对下降，在基线风险1\%时只对应0.2个百分点。政策资源有限时，绝对效应通常比相对效应更接近效用，但不同群体的成本和价值仍需显式权衡。
决策还应考虑估计不确定性和不可逆性。一个平均收益略正但潜在严重损害不确定的政策，与可随时撤回的小成本试点，所需证据阈值不同。因果估计提供后果分布的一部分，决策理论加入效用、风险偏好和学习价值。

\section{从设计到可复核估计}

\subsection{设计原则与正交估计}

观察研究中最重要的工作往往发生在查看结局模型之前。研究者要定义资格、时间零点、处理策略和随访，检查处理组与对照组是否在同一支持区域，并决定哪些变量依据因果结构进入设计。若等看到结果后才调整纳入标准和协变量，模型选择会与期望结论纠缠。

目标试验模拟提供一套组织方式：写出理想随机试验的资格、分配策略、随访、结局、因果对比和分析计划，再说明观察数据如何逐项对应。它不能凭名称消除混杂，但能暴露常见错位。比如新使用者设计避免把已耐受药物的长期使用者与初始未使用者比较；主动对照使就医倾向与适应证更接近；克隆---删失---加权可以比较动态策略，却要求正确处理人工删失和重复单位。

设计阶段的平衡不是最终识别证明。倾向得分匹配后协变量均值接近，只说明已测变量在选定指标上平衡；未测混杂、非线性分布差异和处理版本仍需论证。反过来，极端不平衡是明确警报：目标总体包含数据无法支持的比较，应限制总体、收集新数据或承认外推。

\paragraph{双重稳健与正交化的真实含义}

双重稳健估计结合结果模型$\mu_a(X)$与倾向得分$e(X)$。其一致性通常要求两类模型至少一类正确，而不是任意两个错误会相互抵消。有限样本中，倾向得分接近0或1会放大残差项；结果模型在稀疏区域外推也会不稳定。所谓“双重保险”不能替代重叠诊断和合理变量定义。
正交得分进一步使目标估计对干扰函数的小扰动一阶不敏感。交叉拟合把样本分成若干折，在其他折训练干扰函数，再在留出折计算得分，从而减少复杂学习器对同一噪声过拟合。它允许使用随机森林、提升树或神经网络估计条件均值，但仍要求适当收敛率、独立同分布或相应抽样条件，以及原始识别假设。
实际报告不应只写“使用DML”。应说明得分函数、折数、重复分割、学习器候选、调参准则、截尾规则和标准误计算；比较不同学习器组合与简单参数模型；检查预测性能是否集中在有重叠区域。算法灵活性是估计工具，不是因果结构来源。

\subsection{分层不确定性与可复核分析}

抽样区间只覆盖在目标、识别与模型条件成立时的重复抽样波动。因果研究还存在至少四层不确定性：目标层的不确定性来自处理版本和结局定义；结构层来自DAG、未测混杂与选择机制；估计层来自函数形式、调参与有限样本；迁移层来自目标人群和实施环境变化。把它们压成一个标准误会制造虚假精确。
分层报告可以采用组合方式。主结果给出抽样区间；替代合理调整集和估计器显示建模敏感性；E-value、偏差函数或负对照分析约束未测混杂；权重分布和有效样本量展示重叠；不同人群标准化检验迁移；对处理与结局重新定义展示目标敏感性。各分析应对应明确风险，避免把几十个任意模型称作“稳健性检查”。

决策还需要把不确定性映射到损失。若错过有效治疗与采用有害治疗的代价不同，置信区间跨零并不是唯一标准。可报告达到临床重要阈值的兼容范围、最坏情景净收益和信息价值。统计不确定性由此与行动责任连接，而不是以显著性结束。

\paragraph{从结果表到可复核分析}

可复核报告需要让他人重建从原始问题到最终数字的路径。数据部分说明来源、纳入流程、时间范围、缺失和测量版本；设计部分给出目标试验或DAG、调整理由与支持区域；估计部分给出公式、软件版本、随机种子和诊断；结果部分同时展示绝对与相对尺度、异质性和敏感性；讨论部分区分数据支持与模型承诺。
表格应服务推理，而非堆满系数。基线表用于描述总体与重叠，不以显著性检验选择协变量；权重表报告分位数、截尾和有效样本量；主结果表明确分母、时间窗和效应尺度；敏感性表说明每个情景改变了哪项假设。图形可展示协变量平衡、倾向得分分布、累计风险和异质性校准。
复现还包括失败记录。预先声明哪些诊断会使分析停止或改变目标，例如有效样本量过低、关键层无对照、负对照显示强偏差。若只有成功结果被保存，代码可运行也不等于研究过程可审查。把失败条件写进分析计划，是防止结果导向调整的重要部分。

\section{复杂数据结构与可靠实现}

\subsection{删失、竞争事件与相关数据}

随访数据中的“没有结局”可能表示尚未发生、失访、死亡、退出研究或记录中断。这些状态对应不同目标。若删失概率取决于既往处理和预后，完整案例分析会选择一群更健康或更依从者。逆概率删失加权在给定历史可交换与正值条件下重建未删失总体，但极端权重同样暴露支持缺口。
竞争事件不能一概当删失。研究癌症复发时，死亡使未来复发不再可能；把死亡者删失并假定其仍有潜在复发时间，回答的是一个额外构造的净风险问题。累计发生风险把死亡保留为竞争结局，回答现实世界中的复发概率；复合结局又把死亡与复发按某种规则合并。选择必须服务临床问题，而不是由软件默认决定。

失访敏感性分析可以设定失访者相对已观察者的结局偏移、使用界限或模式混合模型。关键是把无法由数据识别的部分单独参数化。报告失访比例不够，还应展示失访与处理、预后和时间的关系，以及结论需要何种未观察结局分布才会改变。

\paragraph{聚类、网络与相关数据的标准误}

个体记录并不保证独立抽样。学生嵌套学校，患者嵌套医院，居民处于社区和网络；处理也可能在集群层分配。若忽略相关性，点估计有时仍针对同一均值，标准误却通常过小。随机化推断和方差估计应尊重实际分配单位，而不是数据表的最细行。
集群数量少时，常规聚类稳健标准误表现不佳，需要小样本修正、随机化检验或按集群分析。多层模型可以借力估计异质性，但其收缩和分布假设应与因果目标区分。处理在医院层、结局在患者层时，真正有效样本量更接近医院数，而非患者总数。
网络干扰进一步使一个人的处理进入他人结果。标准误问题与目标问题同时改变：需定义曝光映射、网络邻域和策略分配。若网络本身由处理影响，固定网络假设也会失败。设计阶段采用集群或两阶段随机化，通常比事后用复杂相关结构补救更可靠。

\subsection{完整估计实例与数值稳定性}

设目标是某慢病患者在确诊日立即进入强化管理，相对于常规管理，对一年住院风险的差。资格要求连续参保一年且确诊日未住院；处理按确诊后七日内是否进入计划定义；结局为一年内首次住院，死亡作为竞争事件另报。基线共同原因包括年龄、既往住院、合并症、医疗利用和地区资源。
先用DAG说明这些变量阻断计划选择与住院的后门路径，再检查两组倾向得分分布。若极高风险患者全部进入计划，应把目标限制到有比较的患者，而非依赖模型外推。可使用交叉拟合AIPW估计风险差，同时报告结果回归与加权估计；权重截尾应作为敏感性而非隐藏步骤。
诊断包括标准化差异、有效样本量、权重分位数、各医院支持和负对照结局。失访由删失权重处理，并对失访者风险偏移做敏感性分析。最后将结果按地区资源分层，检查从研究医院迁移到基层机构所需的实施稳定性。这个例子显示估计器只是整条链中的一环。

\paragraph{数值稳定性与软件正确性}

正确公式也可能被数值实现破坏。极端权重造成少数单位主导估计，完全分离使倾向模型发散，高维矩阵病态使标准误不稳定，重复交叉拟合产生分割波动。应记录收敛警告、权重最大值、有效样本量和不同随机种子的结果，而不是只保留最终表格。
软件验证可从小型已知答案开始。构造一个有限总体或模拟SCM，手算目标和识别表达式，再检查代码；对权重为1、无缺失、随机处理等特殊情形建立单元测试；比较两个独立实现；确保数据排序、连接和时间窗口没有泄漏。很多“高级方法失败”实际来自单位重复、时间错位或变量编码。
可复现环境应固定依赖版本并保存中间审计摘要，但不必把所有缓存留在项目中。权威源码、配置、数据谱系和最终成品应与可再生成的构建目录分离。清洁目录与严格分析并不冲突，前提是构建命令能够从零恢复全部检查。

\subsection{完整估计案例：从粗比较到目标总体}

设某慢性病管理项目的处理$A$为是否在确诊后七日内进入强化管理，结局$Y$为一年内住院，基线病情$L$简化为低风险和高风险两层。目标总体两层各占一半，但医生更常把高风险患者纳入项目。数据中的条件风险如下：

\begin{center}
\small
\begin{tabularx}{\textwidth}{>{\raggedright\arraybackslash}p{1.35cm}*{4}{>{\centering\arraybackslash}X}}
\toprule
基线层 & 目标占比 & $P(A=1\mid L)$ & $P(Y=1\mid A=1,L)$ & $P(Y=1\mid A=0,L)$\\
\midrule
低风险 & 0.50 & 0.20 & 0.10 & 0.20\\
高风险 & 0.50 & 0.80 & 0.30 & 0.50\\
\bottomrule
\end{tabularx}
\end{center}

在给定$L$后可交换、处理一致且两层都有正值性的假设下，标准化风险为
\begin{align*}
E\{Y(1)\}&=0.5\times0.10+0.5\times0.30=0.20,\\
E\{Y(0)\}&=0.5\times0.20+0.5\times0.50=0.35,
\end{align*}
因而ATE为$-0.15$，即一年住院风险平均下降15个百分点。低风险层的效应为$-0.10$，高风险层为$-0.20$，所以处理分配同时与基线风险和效应异质性有关。

直接比较处理组和对照组却会得到完全不同的印象。由于两个风险层在目标中各占一半，处理组中80\%为高风险，对照组中只有20\%为高风险。于是
\begin{align*}
E(Y\mid A=1)&=0.2\times0.10+0.8\times0.30=0.26,\\
E(Y\mid A=0)&=0.8\times0.20+0.2\times0.50=0.26.
\end{align*}
粗风险差为零，虽然每个风险层内项目都降低住院。这个计算不只展示“混杂可造成偏差”，还说明估计器必须把两组重新带回同一目标总体。

\begin{workedexample}[AIPW的填充与校正]
对处理水平$a=1$，AIPW在每个单位上构造
\[
\widehat m_1(L_i)+\frac{A_i}{\widehat e(L_i)}
\{Y_i-\widehat m_1(L_i)\}.
\]
第一项为所有人填入“如果进入强化管理”的条件风险；第二项只在实际处理者中出现，用逆处理概率放大残差，校正结果模型的系统误差。对$a=0$做同样计算并相减，才得到ATE。若结果模型正确，残差在条件下平均为零；若倾向模型正确，加权残差会修正错误的填充。但如果高风险层从未有对照，任何一项都需外推，“双重稳健”不会创造缺失的比较。
\end{workedexample}

一项完整分析应保留下列次序，而不是从“选什么算法”开始：

\begin{center}
\begin{tikzpicture}[node distance=6mm]
\node[causalbox,text width=2.25cm] (q) {固定问题\\总体与时间};
\node[causalbox,text width=2.25cm,right=of q] (d) {模拟设计\\资格与分配};
\node[causalbox,text width=2.25cm,right=of d] (i) {识别审计\\假设与支持};
\node[causalbox,text width=2.25cm,below=of i] (e) {估计与诊断\\平衡与权重};
\node[causalbox,text width=2.25cm,left=of e] (s) {敏感性\\替代结构};
\node[causalbox,text width=2.25cm,left=of s] (a) {行动解释\\效用与边界};
\draw[causalarrow] (q)--(d);\draw[causalarrow] (d)--(i);\draw[causalarrow] (i)--(e);
\draw[causalarrow] (e)--(s);\draw[causalarrow] (s)--(a);
\end{tikzpicture}
\end{center}

\paragraph{把诊断改写为分析决策}

诊断的价值在于它会改变分析，而不是为附录增加一张图。若倾向得分两组几乎不重叠，预先规定的动作应是收缩目标人群、收集新对照或改用部分识别，不是不断更换分类器。若加权后平衡仍差，应修订处理模型或不再使用该权重。若负对照结果显示强关联，应将未测混杂纳入主结论的不确定性，而不是把它降格为“补充分析”。

同样，估计器之间的差异需要触发解释。在相同目标下，g计算、IPW与AIPW若结果接近，只表示对这些建模选择不敏感；若差异很大，需检查外推、极端权重和函数形式。不能从三个结果中选一个最显著的作为“主模型”。

\begin{failurecheck}[有效样本量与原始行数混为一谈]
假设原始数据有5000人，但少数未处理的高风险患者获得很大权重，使$n_{\mathrm{eff}}$只有450。报告“样本量5000”会让读者误以为每个反事实区域都有充足信息。应同时给出原始行数、独立分配单位数、权重后有效样本量和每个关键层的实际对照数。若信息单位是医院而不是患者，还应把医院数放在推断首位。
\end{failurecheck}

\paragraph{从单一区间到分层不确定性}

对上述住院案例，主表可报告ATE的抽样区间，但还应有三类另行结果。第一类是目标敏感性：把“七日内进入”改为“十四日内进入”，结局和处理含义如何改变。第二类是结构敏感性：未测病情需要多强才会将风险差推回零。第三类是迁移敏感性：项目从大型医院进入基层机构时，依从、容量和对照治疗是否保持。这些问题不宜被合成一个更宽的置信区间，因为它们来自不同证成来源。

\subsection{有限样本推断：随机性来自哪一层}

标准误不是估计器名称的固定配件。它要描述的是，在什么重复过程中，当前估计会如何变动。随机试验中可以把处理分配当作随机性来源；从超总体抽样时可把单位抽样当作来源；固定区县的全部医院数据则未必有一个明确的医院超总体。同一个数字在不同重复概念下可以有不同方差。

\begin{center}
\small
\begin{tabularx}{\textwidth}{>{\bfseries}p{2.4cm}X X}
\toprule
数据生成结构 & 不可拆分的信息单位 & 推断与重抽样原则\\
\midrule
个人随机试验 & 单个受试者，但要保留分层分配 & 按随机化设计或预先指定的稳健方差\\
整群随机试验 & 被随机化的学校、医院或社区 & 在群层重排处理，方差反映群数\\
重复测量 & 同一人的完整时间轨迹 & 按人划分训练折和bootstrap，不拆行\\
匹配或带放回对照 & 匹配组与被重复使用的单位 & 保留匹配及权重依赖，不当作独立两组\\
空间或网络数据 & 由溢出与同一冲击连接的单位 & 由设计定义暴露和分配，不只换用稳健标准误\\
\bottomrule
\end{tabularx}
\end{center}

例如一项学校级干预有30所学校、每校200名学生。若处理在学校层随机分配，6000行记录不能提供6000次独立处理比较。学校内教师、同伴和实施使结局相关；自由度与有效信息更接近30个群。普通个人级bootstrap会在每所学校内抽出大量看似新的独立单位，系统性低估方差。

\paragraph{小样本中稳健方差也会失效}

“使用群稳健标准误”只是开始。群数很少、群规模差异很大或处理群比例失衡时，常用大样本近似仍可能给出过窄区间。可根据设计使用群层随机化检验，使用小样本自由度修正，或者在有限数量匹配群内做精确重排。如果只有两个处理医院和两个对照医院，无论患者数多少，对普遍医院效应的统计推断都很脆弱。

随机化推断也不会自动解决处理后失访、不依从和干扰。它使用已知分配机制评估特定零假设；当结局只在选择后的人中观测，接受的处理版本因人而异，或一人的分配影响他人时，目标和估计方法都要重新定义。

\paragraph{多次交叉拟合与分割不确定性}

在样本不大时，一次随机分折可使某些稀少层几乎只落在一个训练集或评价集，导致DML与AIPW结果随随机种子明显变化。可重复若干次按处理与关键层分层的交叉拟合，报告各次点估计、标准误与最终汇总规则。重复分割不是为了在许多结果中选一个最有利数字，而是暴露训练样本选择对有限样本结果的影响。

需要注意，把多次分割的点估计标准差直接当作总标准误也不正确。每次估计仍共享同一原始样本，分割变动只是算法不稳定的一部分。最终区间应基于目标估计的影响函数或相应理论，并将分割范围作为附加诊断报告。

\begin{failurecheck}[随机拆行制造了虚假外样本]
电子健康记录中同一患者有多次就诊，若按行随机分到训练折和验证折，模型可以通过诊断组合、时间模式甚至隐式标识识别同一人。验证折不再表示新患者。更严重的是，后期就诊包含早期处理及结局的后果，会把处理后信息泄漏到基线预测。分折必须在因果单位与时间零点上完成，并在多中心数据中考虑是按患者、机构还是未来时段检验外推。
\end{failurecheck}

\begin{chapterreview}
识别之后仍可用多种估计器。g计算建模结果，匹配强调设计，IPW建模处理赋值，AIPW与TMLE结合二者，DML以正交得分和交叉拟合容纳高维学习。任何方法都必须与目标量、重叠和研究设计配套；置信区间只覆盖部分不确定性，诊断和敏感性分析不可省略。
\end{chapterreview}

\begin{selfcheck}
\begin{enumerate}
  \item 为什么处理模型的分类准确率不是倾向得分质量的主要标准？
  \item 解释AIPW的两部分分别承担什么功能，“双重稳健”不保证什么？
  \item 交叉拟合与Neyman正交性分别解决哪类机器学习偏差？
  \item 为一项观察研究设计包含重叠、平衡和敏感性分析的报告清单。
\end{enumerate}
\end{selfcheck}

\chapterreadings{
倾向得分的经典论文是\textcite{rosenbaumrubin1983}，应用诊断可读\textcite{austin2011}，匹配估计理论见\textcite{abadieimbens2006}。双重稳健的经典说明见\textcite{bangrobins2005}；潜在结果估计的系统处理见\textcite{imbensrubin2015}；观察性纵向方法见\textcite{hernanrobins2020}；TMLE见\textcite{vanderlaanrose2011}；DML的正式框架见\textcite{chernozhukov2018}。随机试验中的回归调整见\textcite{lin2013}，有限样本和聚类方差问题可结合\textcite{mackinnonwhite1985}与\textcite{cameronmiller2015}阅读。

匹配与分层的历史基础可由\textcite{cochran1968}和\textcite{rubin1973}追溯，现代方法综述见\textcite{stuart2010}；平衡诊断不应被倾向得分模型的预测精度替代，具体比较见\textcite{austin2009}。估计倾向得分的效率结果见\textcite{hirano2003}，直接以平衡为目标的加权和估计可对照\textcite{li2018}与\textcite{imairatkovic2014}。

缺失协变量下的影响函数与加权基础见\textcite{robins1994}，TMLE的原始学习框架见\textcite{vanderlaanrubin2006}。重叠不足时应改变目标或样本而非依赖外推，见\textcite{crump2009}；小样本稳健标准误的有限样本问题见\textcite{imbenskolesar2016}。未测混杂敏感性可分别从非参数界限\parencite{dingvanderweele2016}与回归遗漏变量解释\parencite{cinellihazlett2020}进入。
}
